\chapter{Sistemul de recomandare}
Aplicația practică a modelului antrenat este un sistem de recomandare sub forma unei aplicații web care permite utilizatorului să încarce o
melodie și să primească drept recomandare alte melodii similare.

Când utilizatorul încarcă un fișier audio acesta este trimis sistemului de recomandare aflat pe serverul de backend.
La pornirea serverului este instanțiat un obiect de tipul PredictionSystem. Parametrii cu care este instanțiat
sunt un rezultat al experimentelor realizate în capitolul precedent. Așadar, se specifica dimensiunea inputului
pentru spectrograme Mel ale segmentelor de 15 secunde și o dimensiune a outputului de 24 ce corespunde numărul
final de genuri distincte selectate.

Pipeline-ul de procesare a fișierelor primite este similar cu cel descris în capitolele
\ref{section:extraction} și \ref{section:load}. Fișierul este descărcat local pentru procesare, separat în instrumental și voce,
segmentul audio ce conține vocea este împărțit în segmente de 15 secunde, cu offset-uri
din secundă în secundă, apoi se extrag spectrogramele Mel și se convertesc în \textit{Tensors}.
Datele obținute sunt introduse în modelul rețelei CNN-LSTM pentru a obține vectorii de predicție. Offset-urile mici au fost selectate
deoarece se rezultatele au fost mai bune atunci când existau mai multe instanțe pe care să fie bazată predicția.

Ultimul strat dens conectat are funcția de activare Softmax. Așadar fiecare valoare din vectorii de
predicție rezultați este un procent care reprezintă gradul în care un gen se
regăsește sau nu într-un segment.

\section{Recomandare}

Toate segmentele din setul de date au fost trecute prin model, iar vectorii de predicție au
fost salvați într-un fișier, împreuna cu numele segmentului ce conține genul și
id-ul Spotify al melodiei inițiale. La crearea serverului, obiectul PredictionSystem
instanțiază un obiect KDTree.
Informațiile adiționale menționate sunt încărcate împreună cu toți vectorii de predicție extrași din setul de date.
Aceștia vor fi organizați sub forma unui arbore KDTree pentru a accesa eficient vecinii lor.

Atunci când o melodie este introdusă în sistem,
odată ce au fost obținuți vectorii de predicție pentru segmente, sunt căutați cei mai
apropiați 3 vecini pentru fiecare vector de predicție împreuna cu distanțele dintre
vecini si vectorul de predicție corespunzător. Se selectează cei mai apropiați 12 vecini după
procesul ilustrat în figura \ref{fig:f20}, unde $p_{i}$ este un vector de predicție
pentru un segment al melodiei primite ca input, $i_{ij}$ indexul vectorului de predicție
a unui vecin existent în setul de date, iar $d_{ij}$ este distanța dintre vectorul
de predicție $p_{i}$ și vecinul cu vectorul de predicție la indexul $i_{ij}$.

\begin{figure}[h]
\includegraphics[width=0.9\textwidth]{kdtree.png}
\centering
\captionsetup{justification=centering,font=small,labelfont=small}
\caption{Extragerea vecinilor din KDTree}
\label{fig:f20}

\end{figure}

Lista rezultată în figura \ref{fig:f20} este sortată după distanțele $d_{ij}$ și sunt
returnate id-urile ce corespund celor 12 vecini cu cele mai mici distanțe, identificate prin
indexul $i_{ij}$ al vecinilor. Acești vecini din setul de date sunt, de fapt, melodiile distincte ce vor
fi recomandate, deoarece sunt cele mai similare cu segmentele primite ca input. Acest mod de recomandare
are la bază intuiția că, melodii similare vor fi interpretate similar de rețea și vor avea vectori de predicție cu
valori similare.

\section{Prezicerea genurilor muzicale}
Se amintește faptul că inputurile sunt separate în mai multe segmente. Ca exemplu, pentru un segment de $30$, vom obține $17$ spectrograme Mel ca input pentru rețea.
Având în vedere faptul că pentru o melodie introdusă în sistem nu rezultă un singur vector de predicție, au fost două abordări în prezicerea genurilor muzicale.

Prima abordare a fost calcularea genului muzical pentru fiecare segment în mod clasic, ca fiind genul cu cel mai mare procent din
vectorul de predicție. Astfel, în funcție de numărul de segmente care au fost clasificate ca aparținând unui anumit gen,
s-au calculat procentele pentru fiecare gen prezis.

A doua abordare a fost adunarea tuturor vectorilor de predicție obținuți. Se selectează cele mai mari 3 valori din vectorul sumă obținut,
corespunzătoare genurilor predominante. Restul
genurilor sunt incluse în categoria $others$. Chiar dacă
rezultatele sunt similare cu cele obținute utilizând prima metodă, această
metodă are performanțe mai bune atunci când diferențele dintre procentele din vectorii de predicție obținuți nu sunt mari. Se consideră exemplul din figura \ref{fig:f36}.
Dacă prima metodă ar fi folosită procentele obținute ar fi $66.6\%$ \textit{r-n-b} și $33.3\%$ \textit{pop}. Cu toate acestea, în vectorul sumă se
observă predominanța genului \textit{pop} cu un procent de $40\%$, urmat de \textit{r-n-b} cu $31\%$.

\begin{figure}[h]
\includegraphics[width=0.5\textwidth]{mapping.png}
\centering
\captionsetup{justification=centering,font=small,labelfont=small}
\caption{Exemplu de vectori de predicție și vector sumă}
\label{fig:f36}

\end{figure}

\pagebreak

\section{Experimente}
\label{section:kids}
Un prim experiment a fost compararea genurilor prezise de sistem cu cele reale pentru segmente existente în setul de date. A fost observată
o problemă în cazul genului $kids$. Un număr mare de segmente erau fals prezise ca aparținând acestui gen după cum se observa în figura \ref{fig:f21}.

\begin{figure}[h]
\includegraphics[width=\textwidth]{confusion_better.png}
\centering
\captionsetup{justification=centering,font=small,labelfont=small}
\caption{Genurile prezise inițial de sistem}
\label{fig:f21}

\end{figure}

Ca urmare a acestui rezultat, genul $kids$ a fost eliminat din setul de date, rețeaua a fost re-antrenată și s-a rulat din nou experimentul.
Rezultatele pentru genurile finale selectate sunt ilustrate în figura \ref{fig:f22}. Genul predominant a fost identificat corect în procent de $73.94\%$, iar
în $90.75\%$ dintre cazuri genul real al segmentului se afla printre cele $3$ genuri predominante prezise de sistem.
\pagebreak
\begin{figure}[h]
\includegraphics[width=0.85\textwidth]{confusion_without_kids_better.png}
\centering
\captionsetup{justification=centering,font=small,labelfont=small}
\caption{Genurile prezise de sistem după eliminarea genului kids}
\label{fig:f22}
\end{figure}

Un alt experiment cu scopul de a identifica abilitatea sistemului de a recomanda melodii similare a fost calcularea melodiilor recomandate
care fac parte din același album al unui artist pentru
un segment. Rezultatele pot si observate în figura \ref{fig:f23}.

\begin{figure}[h]
\includegraphics[width=0.85\textwidth]{songs from the same album_without_kids.png}
\centering
\captionsetup{justification=centering,font=small,labelfont=small}
\caption{Numărul de melodii recomandate din acelasi album}
\label{fig:f23}
\end{figure}

Experimentul a fost rulat pentru 119 segmente existente în setul de date. În 80 de cazuri a fost recomandată doar melodia din setul de date din care
face parte segmentul, în 20 de cazuri printre melodiile recomandate se afla melodia inițială și altă melodie din același album, iar în 9 cazuri, mai mult de
2 melodii din același album au fost recomandate.

Segmentele descărcate utilizând Spotify API au doar 30 de secunde. Un subset de 47 de melodii dintre segmentele utilizate pentru experimentul anterior au fost
descărcate în întregime, iar experimentul a fost rulat din nou. Rezultatele pot fi observate în figurile \ref{fig:f24} și \ref{fig:f25}. Genul predominant a fost identificat corect în procent de $57.44\%$, iar
în $80.85\%$ dintre cazuri genul real al segmentului se afla printre cele $3$ genuri predominante prezise de sistem.

\begin{figure}[h]
\centering
\includegraphics[width=0.9\textwidth]{confusion_full_songs_better_without_kids.png}
\captionsetup{justification=centering,font=small,labelfont=small}
\caption{Genurile prezise de sistem pentru melodiile întregi}
\label{fig:f24}
\end{figure}

\begin{figure}[h]
\centering
\includegraphics[width=0.8\textwidth]{songs_from_the_same_album_full_songs_without_kids.png}
\captionsetup{justification=centering,font=small,labelfont=small}
\caption{Numărul de melodii recomandate din acelasi album pentru melodiile întregi}
\label{fig:f25}
\end{figure}

În 31 dintre cazuri, segmentul existent în setul de date din melodia întreagă nu a fost identificat. Acest lucru poate fi datorat offset-ului utilizat de Spotify
pentru selectarea segmentului, segmentele nu se suprapun perfect peste cel existent in setul de date deci se găsesc alte melodii cu segmente mai
similare decât cea originală.
